\documentclass{article}%
\usepackage[T1]{fontenc}%
\usepackage[utf8]{inputenc}%
\usepackage{lmodern}%
\usepackage{textcomp}%
\usepackage{lastpage}%
\usepackage{geometry}%
\geometry{margin=2.2cm}%
%
\usepackage{amsmath}%
\usepackage{amssymb}%
\usepackage{booktabs}%
\usepackage{longtable}%
\usepackage{array}%
\usepackage{hyperref}%
\hypersetup{hidelinks}%
\title{IME 2027 {-} Matemática (claude\_sonnet\_5\_5) x Poliedro\\[4pt]\large Comparação com o gabarito Poliedro\thanks{Gerado com gpt-resolve: \url{https://github.com/lgabs/gpt-resolve}}}%
\date{09/11/2026}%
%
\begin{document}%
\normalsize%
\maketitle%
\section*{Resumo}\textbf{Nota média: 10,0/10} (12 de 12 itens com a resolução correta; cada item vale o mesmo dentro da questão).\par\medskip%
\begin{longtable}{@{}c c p{0.78\textwidth}@{}}\toprule Questão & Nota & Resumo \\ \midrule
1 & 10,0 & A resolução e o gabarito Poliedro concordam: a área de P ? I é 12$\surd$3 - 18 $\approx$ 2,785. Confirmei o valor recortando os polígonos numericamente: a interseção é um dodecágono de área 18 - 9$\surd$3 e cada hexágono tem área 3$\surd$3/2. A resolução está correta e bem justificada. Ela mostra que nenhuma das 12 retas é redundante e confere o resultado por um segundo caminho, somando a interseção às 12 pontas da estrela. Não há divergência a levar ao professor; só uma ressalva de redação, de pouca importância, sobre a convenção de numeração das raízes. \\
2 & 10,0 & A resolução está correta e coincide com o gabarito Poliedro: P = 2/5. As raízes são cis(k$\pi$/3), k=1..5. Usando |z|=1, a inequação se reduz a Re z $\geq$ 1/2, e só z = e\^{}\{$\pm$i$\pi$/3\} a satisfazem, com igualdade. Conferi as contas numericamente. Não há divergências nem erros. Os únicos pontos de atenção são a hipótese implícita de equiprobabilidade e o fato de a resposta depender do sinal $\leq$ na fronteira. \\
3 & 10,0 & A resolução e o gabarito chegam ao mesmo resultado, x = 7, e os dois estão certos. Refiz a conta, f(x)+f(-1/x) = 1 - 3\^{}\{1-x\}/2 para x $\geq$ 1, e conferi com frações exatas em Python: o primeiro inteiro que satisfaz a desigualdade é 7. A resolução está mais completa que o gabarito. Ela justifica que inteiros negativos dão soma 0,5 e que x=0 fica fora do domínio, coisa que o Poliedro só afirma sem justificar. \\
4 & 10,0 & A resolução e o gabarito Poliedro dão o mesmo valor, 1345, e o valor está correto. Os dois usam o mesmo argumento: $\Pi$(y\_i+x) = x$^3${-}4x+2, depois a redução módulo x$^2${-}7x{-}3 e as relações de Girard. Confirmei o resultado recalculando numericamente as raízes em Python. As contas e as checagens da resolução (sinal, discriminantes) estão certas, e não há pontos para levar ao professor. \\
5 & 10,0 & A resolução está correta e não contradiz o gabarito do Poliedro. Ela resolve o problema como foi impresso (5540): a solução é única, S $\approx$ 11196,5, e conferi isso numericamente. Também identifica o provável erro de digitação, pois com 5440 obtém S = 10912, que é a resposta do gabarito. O gabarito acerta ao supor o erro e ao chegar a 10912. Só exagera ao dizer que não há 'forma fechada' com 5540: Cardano dá a raiz por radicais, apenas sem valor simples. Como o PDF original de fato traz 5540, a questão pode ser anulada ou ter mais de uma resposta aceita. \\
6 & 10,0 & A resolução e o gabarito Poliedro estão de acordo: P = 1/3, e ambos estão certos. A transcrição do enunciado bate com a prova original (página 2, 6ª questão). A solução está completa e bem justificada. Ela mostra que o número de bijeções é 8! para cada comprimento k do ciclo que contém o 1, confere que a soma dá 9! e soma os casos k = 1, 3 e 9. Confirmei o resultado enumerando todas as 9! bijeções em Python: 120960 são favoráveis. \\
7 & 10,0 & A resolução está correta e completa. Ela chega a tr(AB) = C(m+n+2, n+1) = C(2026, 1001), a mesma expressão do gabarito Poliedro. A divisibilidade por 1013 está provada com rigor: 1013 é primo, 1001$\cdot$C(2026,1001) = 2$\cdot$1013$\cdot$C(2025,1000) e mdc(1001, 1013) = 1. Conferi tudo em Python, tanto a fórmula do traço quanto C(2026,1001) $\equiv$ 0 (mod 1013). O gabarito tem a conclusão certa, mas a justificativa tem um passo falso: '1013 não divide 1025!$\cdot$1001!', quando na verdade 1013 | 1025!. A divisibilidade vale porque o expoente de 1013 é 2 em 2026! e 1 em 1025!$\cdot$1001!. Veredito: as respostas coincidem e a resolução merece nota integral. \\
8 & 10,0 & A resolução está correta e completa, e bate com o gabarito Poliedro (mesmo método: prostaférese, diferença de tangentes e soma telescópica). Ela justifica explicitamente que os denominadores são não nulos a partir das hipóteses. A identidade também foi confirmada numericamente. Não há divergência: as duas estão certas. \\
9 & 10,0 & A resolução está correta e bate com o gabarito Poliedro. As duas chegam ao lugar geométrico 4x$^2$/49 + y$^2$ = 1 e aos quatro pontos ($\pm$2$\surd$3, $\pm$1/7). Conferi o enunciado no PDF original e refiz as contas, inclusive numericamente. A resolução é até mais rigorosa que o gabarito, porque exclui os casos degenerados P = ($\pm$4, 0), sem efeito sobre a resposta. Não há erro a apontar. \\
10 & 10,0 & A resolução e o gabarito Poliedro dão a mesma resposta nos dois itens, e ambos estão corretos: (a) V\_max = 192$\surd$3 cm$^3$ $\approx$ 332,55 cm$^3$; (b) o prisma é um cubo com altura e lados da base iguais a 4$\surd$3 cm. Refiz as contas e fiz uma busca numérica, e os valores se confirmam. A resolução é mais completa que a do gabarito. Ela trata um quadrilátero inscrito qualquer (área $\leq$ 2r$^2$, com igualdade só no quadrado) e deduz que h $\leq$ 2d a partir da condição de estar contido na esfera. Também mostra que a cota 192$\surd$3 vale para qualquer prisma reto quadrangular dentro da esfera, inclusive quando a face com vértices na esfera é lateral. O gabarito supõe uma base retangular inscrita com os 8 vértices na esfera, sem justificar, mas isso não altera o resultado. As ressalvas restantes são pequenas e de redação. \\
\bottomrule\end{longtable}%
\section*{Principais pontos de dúvida}%
\begin{itemize}%
\item%
\textbf{Q5} \textit{(statement)} No PDF original está impresso x2+x4+x6+x8 = 5540. Com esse valor a resposta é irracional ($\approx$11196,5). Tudo indica erro da banca no lugar de 5440, que dá q = 2, x\_k = 2\^{}(k+4) e S = 10912; o produto 2\^{}45 bate exatamente com q = 2. A nota depende de como a banca tratar o problema (anular, aceitar 10912 ou aceitar $\approx$11196,5). A resolução cobre os dois casos.%
\item%
\textbf{Q7} \textit{(key)} A resolução do gabarito Poliedro afirma que 'o primo 1013 não divide 1025!$\cdot$1001!'. Isso é falso: como 1013 $\leq$ 1025, 1013 é fator de 1025! (conferido: 1025! $\equiv$ 0 mod 1013). O argumento certo usa expoentes (Legendre): v\_1013(2026!) = 2, v\_1013(1025!) = 1, v\_1013(1001!) = 0, logo v\_1013(C(2026,1001)) = 1. Outro caminho é o da resolução: 1001$\cdot$C(2026,1001) = 2$\cdot$1013$\cdot$C(2025,1000) com mdc(1001, 1013) = 1. A conclusão continua certa, mas uma correção que copie o passo do gabarito estaria demonstrando com um argumento inválido. Isso pode pesar numa questão de 'mostre que'.%
\end{itemize}%
\section*{Questão 1}%
\href{https://poliedroresolve.sistemapoliedro.com.br/vestibulares/ime/2027/ime-2-fase-dia-1-26-10-2026-2-fase-dia-1/dissertativa-1-matematica-2-fase-dia-1-ime-2027}{Resolução do gabarito}\par\smallskip%
\begin{longtable}{@{}p{0.06\textwidth} p{0.33\textwidth} p{0.33\textwidth} p{0.1\textwidth} p{0.08\textwidth}@{}}\toprule Item & Resolução & Gabarito & Veredito & Certo \\ \midrule
Área de P ? I & 12$\surd$3 - 18 $\approx$ 2,785 & 12$\surd$3 - 18 ($\approx$ 2,785) & igual & ambos \\
\bottomrule\end{longtable}%
\begin{itemize}%
\item%
\textbf{Área de P ? I)} A resolução e o gabarito dão o mesmo valor e usam o mesmo caminho: cada hexágono tem área 3$\surd$3/2, a interseção é um dodecágono regular de apótema $\surd$3/2 com área 12a$^2$$\cdot$tg15$^\circ$ = 18 - 9$\surd$3, e a inclusão--exclusão dá 3$\surd$3 - (18 - 9$\surd$3) = 12$\surd$3 - 18. Refiz a conta em Python recortando um hexágono pelo outro (Sutherland--Hodgman). Saiu uma interseção de 12 vértices com área 2,41154 = 18 - 9$\surd$3 e uma união de 2,78461 = 12$\surd$3 - 18. Também conferi a verificação pelas pontas da estrela: T = $\frac12$(2$\surd$3 - 3)(1 - $\surd$3/2) = 7$\surd$3/4 - 3 $\approx$ 0,0311, e A\_D + 12T = 12$\surd$3 - 18. A cota superior 3 é a área do dodecágono regular inscrito de raio 1, que é a envoltória convexa da união; essa cota também está certa.%
\end{itemize}%
\paragraph{Pontos de dúvida}%
\begin{itemize}%
\item%
\textit{(solution)} A resolução diz que o resultado não depende de onde começa a numeração, porque um hexágono regular com vértices entre as 12 raízes só pode ser uma das duas classes alternadas. A afirmação está certa e é coerente com o enunciado, que diz que P e I são regulares. Mesmo assim, ela é dada sem demonstração. Também fica implícita a numeração usual z\_k = cis(k$\pi$/6), que é a mesma do gabarito.%
\end{itemize}%
\section*{Questão 2}%
\href{https://poliedroresolve.sistemapoliedro.com.br/vestibulares/ime/2027/ime-2-fase-dia-1-26-10-2026-2-fase-dia-1/dissertativa-2-matematica-2-fase-dia-1-ime-2027}{Resolução do gabarito}\par\smallskip%
\begin{longtable}{@{}p{0.06\textwidth} p{0.33\textwidth} p{0.33\textwidth} p{0.1\textwidth} p{0.08\textwidth}@{}}\toprule Item & Resolução & Gabarito & Veredito & Certo \\ \midrule
Q2 & P = 2/5 (40\%) & P = 2/5 & igual & ambos \\
\bottomrule\end{longtable}%
\begin{itemize}%
\item%
\textbf{Q2)} Conferi o enunciado no PDF original e ele bate com a transcrição. A equação é 1+z+...+z\^{}5 = (z\^{}6-1)/(z-1) = 0, e suas raízes são as 5 raízes sextas da unidade diferentes de 1: cis(k$\pi$/3), k=1..5. Todas são simples e distintas. Como |z|=1, vale |z-3/2|$^2$ = 13/4 - 3Re z, e a inequação equivale a Re z $\geq$ 1/2. Conferi numericamente: |z-3/2|$^2$ vale 7/4 para k=1 e 5, 19/4 para k=2 e 4, e 25/4 para k=3, o que coincide com os módulos $\surd$7/2, $\surd$19/2 e 5/2 do gabarito. Só z = e\^{}\{$\pm$i$\pi$/3\} satisfazem, e com igualdade, que o sinal $\leq$ permite. Logo P = 2/5. A resolução e o gabarito chegam ao mesmo resultado pelo mesmo raciocínio. A resolução ainda acrescenta a interpretação geométrica e a observação correta de que z=1 não é raiz.%
\end{itemize}%
\paragraph{Pontos de dúvida}%
\begin{itemize}%
\item%
\textit{(statement)} O enunciado não diz explicitamente que as raízes são equiprováveis. A resolução e o gabarito adotam a interpretação natural: escolha uniforme entre as 5 raízes distintas, que são todas simples, então não há multiplicidade a ponderar.%
\item%
\textit{(statement)} As duas raízes favoráveis ficam exatamente sobre a fronteira do disco (|z-3/2| = $\surd$7/2). A resposta 2/5 depende de o sinal ser $\leq$. Com desigualdade estrita, a probabilidade seria 0. A resolução destaca esse ponto corretamente.%
\end{itemize}%
\section*{Questão 3}%
\href{https://poliedroresolve.sistemapoliedro.com.br/vestibulares/ime/2027/ime-2-fase-dia-1-26-10-2026-2-fase-dia-1/dissertativa-3-matematica-2-fase-dia-1-ime-2027}{Resolução do gabarito}\par\smallskip%
\begin{longtable}{@{}p{0.06\textwidth} p{0.33\textwidth} p{0.33\textwidth} p{0.1\textwidth} p{0.08\textwidth}@{}}\toprule Item & Resolução & Gabarito & Veredito & Certo \\ \midrule
Questão 3 (menor inteiro x com f(x)+f(-1/x) $\geq$ 0,999) & x = 7 & x = 7 & igual & ambos \\
\bottomrule\end{longtable}%
\begin{itemize}%
\item%
\textbf{Questão 3 (menor inteiro x com f(x)+f(-1/x) $\geq$ 0,999))} Refiz a questão e a resposta é 7. A condição g\_n(x)=1 equivale a x > -1/n. Para x $\geq$ 0 isso vale para todo n, então f(x)=1/2. Para x inteiro $\geq$ 1, g\_n(-1/x)=1 só quando n < x; assim f(-1/x)=(1-3\^{}\{1-x\})/2 e a soma dá 1 - 3\^{}\{1-x\}/2. A desigualdade vira 3\^{}\{x-1\} $\geq$ 500, logo x-1 $\geq$ 6 e x $\geq$ 7. Valores: h(6)=1-1/486$\approx$0,99794 (não serve) e h(7)=1-1/1458$\approx$0,99931 (serve). Para x inteiro $\leq$ -1: f(x)=0 e f(-1/x)=1/2, então a soma é 0,5. Em x=0 a expressão f(-1/x) não existe. Conferi com frações exatas em Python para x de -5 a 9, e o primeiro inteiro que satisfaz é 7. Conferi também o enunciado transcrito com o PDF e está fiel.%
\end{itemize}%
\paragraph{Pontos de dúvida}%
\begin{itemize}%
\item%
\textit{(key)} O gabarito do Poliedro diz que x $\in$ $\mathbb{Z}$+ sem justificar. Como a pergunta é o MENOR inteiro, falta mostrar que nenhum inteiro negativo serve. Para x $\leq$ -1, f(x)=0 e f(-1/x)=1/2, e a soma 0,5 não atinge 0,999. Também falta dizer que x=0 fica fora porque -1/x não existe. A resolução conferida faz essa análise; o gabarito pula. Isso pode pesar numa correção rigorosa da justificativa, mas não muda a resposta.%
\item%
\textit{(solution)} No caso x=n, a resolução usa que g\_x(-1/x)=0, porque a desigualdade é estrita (-1/x < -1/n). Isso decide se o último termo entra na soma, e ela trata corretamente. Se essa fronteira fosse tratada errado, a resposta mudaria para x=6.%
\end{itemize}%
\section*{Questão 4}%
\href{https://poliedroresolve.sistemapoliedro.com.br/vestibulares/ime/2027/ime-2-fase-dia-1-26-10-2026-2-fase-dia-1/dissertativa-4-matematica-2-fase-dia-1-ime-2027}{Resolução do gabarito}\par\smallskip%
\begin{longtable}{@{}p{0.06\textwidth} p{0.33\textwidth} p{0.33\textwidth} p{0.1\textwidth} p{0.08\textwidth}@{}}\toprule Item & Resolução & Gabarito & Veredito & Certo \\ \midrule
Q4 (valor do produto) & 1345 & 1345 & igual & ambos \\
\bottomrule\end{longtable}%
\begin{itemize}%
\item%
\textbf{Q4 (valor do produto))} Conferi o enunciado no PDF original, e a transcrição está fiel. Os dois usam o mesmo método. Tomando y = {-}x em y$^3${-}4y{-}2 = $\Pi$(y{-}y\_i), obtém{-}se $\Pi$(y\_i+x) = x$^3${-}4x+2. Pela equação quadrática, x\_j$^3$ = 52x\_j+21, o que reduz cada fator a 48x\_j+23. Por Girard (soma = 7, produto = {-}3): 2304$\cdot$({-}3) + 1104$\cdot$7 + 529 = 1345. O gabarito chega a ({-}48x\_j{-}23); como são dois fatores, os sinais negativos se cancelam e o produto é o mesmo. Recalculei numericamente em Python as raízes y $\approx$ {-}1,6751; {-}0,5391; 2,2143 e x $\approx$ 7,4051; {-}0,4051, e o produto dos seis fatores deu 1345,0000. O discriminante da cúbica, 256{-}108 = 148, também está correto.%
\end{itemize}%
\section*{Questão 5}%
\href{https://poliedroresolve.sistemapoliedro.com.br/vestibulares/ime/2027/ime-2-fase-dia-1-26-10-2026-2-fase-dia-1/dissertativa-5-matematica-2-fase-dia-1-ime-2027}{Resolução do gabarito}\par\smallskip%
\begin{longtable}{@{}p{0.06\textwidth} p{0.33\textwidth} p{0.33\textwidth} p{0.1\textwidth} p{0.08\textwidth}@{}}\toprule Item & Resolução & Gabarito & Veredito & Certo \\ \midrule
Questão 5 --- x1+x3+x5+x7+x9 & Com o dado impresso (5540): S = 512(t$^{4}$-3t$^2$+1) = -512t$^2$+5540t+512 $\approx$ 11196,5, onde t = q+1/q $\approx$ 2,5116 é a única raiz t>2 de 128t$^3$-256t-1385=0 (x5 = 512). Aponta 5540 como provável erro de digitação de 5440; com 5440: q=2, x\_k = 2\^{}(k+4) e S = 10912. & Com 5540 a questão não tem solução em forma fechada. Supondo erro de digitação (5440): q=2 e S = 10912. Sugere anular a questão. & equivalente & ambos \\
\bottomrule\end{longtable}%
\begin{itemize}%
\item%
\textbf{Questão 5 --- x1+x3+x5+x7+x9)} As duas leituras não se contradizem: ambas chegam a x5 = 512 e a 512(u$^3$-2u) = 5540 com u = q+1/q, e ambas dão 10912 com o valor 5440. Conferi no PDF original que o número impresso é mesmo 5540, então a transcrição está fiel. Refiz as contas: (i) o teste das raízes racionais de 128t$^3$-256t-1385 não dá nenhuma raiz racional; (ii) a raiz real por bissecção é t $\approx$ 2,511600, o que dá q $\approx$ 2,015427, r $\approx$ 1,0111 > 0 e a $\approx$ 4,956. Com esses valores a soma dos pares dá 5540,000 e o produto dos ímpares dá 2\^{}45, como pedido. (iii) S = $\Sigma$ ímpares $\approx$ 11196,4994, e as três expressões da resolução coincidem: soma direta, 512(t$^{4}$-3t$^2$+1) e -512t$^2$+5540t+512. (iv) Com x\_k = 2\^{}(k+4), a soma dos pares é 5440 e a dos ímpares é 10912. A unicidade também está certa: g(t) = t$^3$-2t é crescente para t>2 e g(2) = 4 < 1385/128. Do ponto de vista matemático, a resolução é mais completa: ela resolve o problema como foi impresso (que tem solução real única, só que feia) e também a versão provavelmente pretendida. O gabarito diz que não há 'forma fechada', o que não é bem verdade: o discriminante é negativo, a cúbica tem uma única raiz real e a fórmula de Cardano a expressa por radicais. O gabarito acerta, porém, que não existe valor 'bonito'.%
\end{itemize}%
\paragraph{Pontos de dúvida}%
\begin{itemize}%
\item%
\textit{(statement)} No PDF original está impresso x2+x4+x6+x8 = 5540. Com esse valor a resposta é irracional ($\approx$11196,5). Tudo indica erro da banca no lugar de 5440, que dá q = 2, x\_k = 2\^{}(k+4) e S = 10912; o produto 2\^{}45 bate exatamente com q = 2. A nota depende de como a banca tratar o problema (anular, aceitar 10912 ou aceitar $\approx$11196,5). A resolução cobre os dois casos.%
\item%
\textit{(key)} O gabarito diz que, com 5540, 'não tem solução em forma fechada'. Isso é impreciso. O problema tem solução real única (t $\approx$ 2,5116, S $\approx$ 11196,5), e a fórmula de Cardano dá essa raiz por radicais: t = ?(1385/256 + $\surd$((1385/256)$^2$ - 8/27)) + ?(1385/256 - $\surd$((1385/256)$^2$ - 8/27)). O que não existe é um valor racional ou simples.%
\item%
\textit{(solution)} Para o dado impresso, a resolução declara como resposta principal um valor aproximado ($\approx$11196,5), escrito como S = 512(t$^{4}$-3t$^2$+1) com t raiz da cúbica. Não dá a expressão exata por Cardano. Isso é aceitável, mas um corretor rigoroso poderia pedir a forma exata.%
\end{itemize}%
\section*{Questão 6}%
\href{https://poliedroresolve.sistemapoliedro.com.br/vestibulares/ime/2027/ime-2-fase-dia-1-26-10-2026-2-fase-dia-1/dissertativa-6-matematica-2-fase-dia-1-ime-2027}{Resolução do gabarito}\par\smallskip%
\begin{longtable}{@{}p{0.06\textwidth} p{0.33\textwidth} p{0.33\textwidth} p{0.1\textwidth} p{0.08\textwidth}@{}}\toprule Item & Resolução & Gabarito & Veredito & Certo \\ \midrule
Q6 & P = 1/3 & P = 1/3 & igual & ambos \\
\bottomrule\end{longtable}%
\begin{itemize}%
\item%
\textbf{Q6)} As duas respostas são iguais e o raciocínio é o mesmo: f\^{}(9)(1)=1 vale exatamente quando o ciclo que contém o 1 tem comprimento k divisor de 9, ou seja, k $\in$ \{1, 3, 9\}. Para cada k, o número de bijeções é C(8,k-1)$\cdot$(k-1)!$\cdot$(9-k)! = 8!. O gabarito confere o caso k=3: C(8,2)$\cdot$2$\cdot$6! = 40320 = 8!. Então P = 3$\cdot$8!/9! = 1/3. Rodei uma enumeração completa das 9! = 362880 bijeções em Python e encontrei 120960 favoráveis (razão 1/3), o que confirma o valor.%
\end{itemize}%
\section*{Questão 7}%
\href{https://poliedroresolve.sistemapoliedro.com.br/vestibulares/ime/2027/ime-2-fase-dia-1-26-10-2026-2-fase-dia-1/dissertativa-7-matematica-2-fase-dia-1-ime-2027}{Resolução do gabarito}\par\smallskip%
\begin{longtable}{@{}p{0.06\textwidth} p{0.33\textwidth} p{0.33\textwidth} p{0.1\textwidth} p{0.08\textwidth}@{}}\toprule Item & Resolução & Gabarito & Veredito & Certo \\ \midrule
7 & tr(AB) = C(m+n+2, n+1) = C(2026, 1001). É divisível pelo primo 1013 porque 1001$\cdot$C(2026,1001) = 2026$\cdot$C(2025,1000) = 2$\cdot$1013$\cdot$C(2025,1000) e mdc(1013, 1001) = 1 (lema de Euclides). & tr(AB) = C(m+n+2, n+1) = C(2026, 1001) = 2026!/(1025!$\cdot$1001!), divisível por 1013. & igual & ambos \\
\bottomrule\end{longtable}%
\begin{itemize}%
\item%
\textbf{7)} A conclusão e a fórmula fechada são iguais nas duas. Refiz as contas. Os três elementos da diagonal saem por Vandermonde: c11 = C(m+n, n-1) (k = j-3, de -2 a n-1, completo), c22 = C(m+n, n) e c33 = C(m+n+1, n+1). Aplicando Stifel duas vezes, tr = C(m+n+2, n+1). Para m = 1024 e n = 1000 isso dá C(2026, 1001). Conferi em Python a fórmula do traço, somando as entradas diretamente para todos os pares com 1 $\leq$ n < m $\leq$ 24, e também o caso real (1024, 1000): C(2026,1001) mod 1013 = 0. Confirmei ainda que 1013 é primo e que os restos listados na resolução (1,2,3,5,1,12,10,6,1,27,21) estão corretos. O argumento de divisibilidade da resolução é rigoroso. O do gabarito tem um passo falso: diz que '1013 não divide 1025!$\cdot$1001!', mas 1013 $\leq$ 1025, logo 1013 | 1025!. O argumento certo conta expoentes: v\_1013(2026!) = 2 (os fatores 1013 e 2026), v\_1013(1025!) = 1 e v\_1013(1001!) = 0, então v\_1013(tr) = 1 $\geq$ 1. A conclusão do gabarito é verdadeira, mas a justificativa publicada não vale. Por isso o item fica como 'both' quanto à resposta final, e a falha do gabarito vai nas dúvidas.%
\end{itemize}%
\paragraph{Pontos de dúvida}%
\begin{itemize}%
\item%
\textit{(key)} A resolução do gabarito Poliedro afirma que 'o primo 1013 não divide 1025!$\cdot$1001!'. Isso é falso: como 1013 $\leq$ 1025, 1013 é fator de 1025! (conferido: 1025! $\equiv$ 0 mod 1013). O argumento certo usa expoentes (Legendre): v\_1013(2026!) = 2, v\_1013(1025!) = 1, v\_1013(1001!) = 0, logo v\_1013(C(2026,1001)) = 1. Outro caminho é o da resolução: 1001$\cdot$C(2026,1001) = 2$\cdot$1013$\cdot$C(2025,1000) com mdc(1001, 1013) = 1. A conclusão continua certa, mas uma correção que copie o passo do gabarito estaria demonstrando com um argumento inválido. Isso pode pesar numa questão de 'mostre que'.%
\item%
\textit{(solution)} A resolução diz que conferiu a fórmula numericamente em alguns pares. Isso não faz parte da demonstração e não atrapalha. A prova formal (Vandermonde com o cuidado nos limites, Stifel, primalidade de 1013 e o lema de Euclides) está completa e confere com o enunciado original.%
\end{itemize}%
\section*{Questão 8}%
\href{https://poliedroresolve.sistemapoliedro.com.br/vestibulares/ime/2027/ime-2-fase-dia-1-26-10-2026-2-fase-dia-1/dissertativa-8-matematica-2-fase-dia-1-ime-2027}{Resolução do gabarito}\par\smallskip%
\begin{longtable}{@{}p{0.06\textwidth} p{0.33\textwidth} p{0.33\textwidth} p{0.1\textwidth} p{0.08\textwidth}@{}}\toprule Item & Resolução & Gabarito & Veredito & Certo \\ \midrule
Q8 (demonstração) & Identidade demonstrada. Usa a prostaférese cos $\alpha$ + cos((2k+1)$\alpha$) = 2cos(k$\alpha$)cos((k+1)$\alpha$) e escreve cada termo como {[}tg((k+1)$\alpha$) - tg(k$\alpha$){]}/(2 sen $\alpha$). A soma telescopa para {[}tg((n+1)$\alpha$) - tg $\alpha${]}/(2 sen $\alpha$) = sen(n$\alpha$)/(sen 2$\alpha$ $\cdot$ cos((n+1)$\alpha$)). & Demonstração: S\_n = sen(n$\alpha$)/(sen(2$\alpha$)$\cdot$cos((n+1)$\alpha$)), pelo mesmo caminho (prostaférese, diferença de tangentes e soma telescópica). & igual & ambos \\
\bottomrule\end{longtable}%
\begin{itemize}%
\item%
\textbf{Q8 (demonstração))} A resolução segue o mesmo caminho do gabarito e chega à mesma conclusão. Conferi cada passo. (1) A fórmula de prostaférese com p=(2k+1)$\alpha$ e q=$\alpha$ dá (p+q)/2=(k+1)$\alpha$ e (p-q)/2=k$\alpha$, então está correta. (2) tg A - tg B = sen(A-B)/(cos A cos B), com A-B=$\alpha$ no termo geral e A-B=n$\alpha$ no fechamento. (3) A soma telescopa corretamente. (4) 2 sen $\alpha$ cos $\alpha$ = sen 2$\alpha$. A resolução também justifica bem que os denominadores não se anulam: cos(j$\alpha$)$\neq$0 para j=1..n+1 e sen 2$\alpha$$\neq$0 (porque sen $\alpha$$\neq$0 e cos $\alpha$$\neq$0, caso j=1). Nisso ela é até mais cuidadosa que o gabarito resumido. Uma checagem numérica em python, com $\alpha$ e n aleatórios, confirmou a identidade (erro relativo de 1e{-}13 a 1e{-}16). O enunciado transcrito confere com o PDF original: lá aparece cos(2k+1)$\alpha$, sem parênteses externos, que significa o mesmo.%
\end{itemize}%
\paragraph{Pontos de dúvida}%
\begin{itemize}%
\item%
\textit{(solution)} A verificação 'n=0 (soma vazia)' só vale se $\mathbb{N}$ incluir o 0. Isso é inofensivo: o caso trivial dá 0 = 0 e não faz parte da prova principal. Já a frase 'uma checagem numérica com valores aleatórios confirmou a igualdade' não tem valor de demonstração numa prova discursiva. Também é inofensiva, porque a demonstração algébrica está completa.%
\end{itemize}%
\section*{Questão 9}%
\href{https://poliedroresolve.sistemapoliedro.com.br/vestibulares/ime/2027/ime-2-fase-dia-1-26-10-2026-2-fase-dia-1/dissertativa-9-matematica-2-fase-dia-1-ime-2027}{Resolução do gabarito}\par\smallskip%
\begin{longtable}{@{}p{0.06\textwidth} p{0.33\textwidth} p{0.33\textwidth} p{0.1\textwidth} p{0.08\textwidth}@{}}\toprule Item & Resolução & Gabarito & Veredito & Certo \\ \midrule
Pontos de interseção do lugar geométrico de M com x = $\pm$c & (2$\surd$3, $\pm$1/7) e (-2$\surd$3, $\pm$1/7), quatro pontos no total & (2$\surd$3, 1/7), (2$\surd$3, -1/7), (-2$\surd$3, 1/7), (-2$\surd$3, -1/7) & igual & ambos \\
Lugar geométrico de M (resultado intermediário) & 4x$^2$/49 + y$^2$ = 1, sem os pontos com y = 0 & 4x$^2$/49 + y$^2$ = 1 & equivalente & ambos \\
\bottomrule\end{longtable}%
\begin{itemize}%
\item%
\textbf{Pontos de interseção do lugar geométrico de M com x = $\pm$c)} Refiz a conta: c$^2$ = 16 - 4 = 12, logo c = 2$\surd$3. A normal em P = (x0, y0) tem direção (x0, 4y0) e corta Y = 0 em t = -1/4, o que dá Q = (3x0/4, 0) e M = (7x0/8, y0/2). Com x$^2$ = 12 no lugar 4x$^2$/49 + y$^2$ = 1, sai y$^2$ = 1/49 e y = $\pm$1/7. Como 2$\surd$3 $\approx$ 3,464 < 7/2, as retas de fato cortam o lugar. Também conferi numericamente para vários pontos P da elipse.%
\item%
\textbf{Lugar geométrico de M (resultado intermediário))} É a mesma equação. A resolução apenas retira os pontos ($\pm$7/2, 0), que viriam de P = ($\pm$4, 0). Nesses pontos a normal é o próprio eixo x e Q não fica definido. Essa exclusão é mais rigorosa e não muda os pontos pedidos, que têm y = $\pm$1/7 $\neq$ 0.%
\end{itemize}%
\paragraph{Pontos de dúvida}%
\begin{itemize}%
\item%
\textit{(key)} O gabarito não trata os casos degenerados P = ($\pm$4, 0), em que a normal coincide com o eixo x e Q fica indeterminado. Por isso, ele apresenta o lugar como a elipse inteira, incluindo ($\pm$7/2, 0). Isso não muda a resposta pedida. A resolução trata esses casos corretamente.%
\item%
\textit{(key)} O gabarito escreve o coeficiente angular da normal como 4y0/x0, o que só vale para x0 $\neq$ 0. O resultado Q = (3x0/4, 0) continua valendo em x0 = 0, onde a normal é o eixo y, como a resolução observa. É só uma questão de redação.%
\end{itemize}%
\section*{Questão 10}%
\href{https://poliedroresolve.sistemapoliedro.com.br/vestibulares/ime/2027/ime-2-fase-dia-1-26-10-2026-2-fase-dia-1/dissertativa-10-matematica-2-fase-dia-1-ime-2027}{Resolução do gabarito}\par\smallskip%
\begin{longtable}{@{}p{0.06\textwidth} p{0.33\textwidth} p{0.33\textwidth} p{0.1\textwidth} p{0.08\textwidth}@{}}\toprule Item & Resolução & Gabarito & Veredito & Certo \\ \midrule
a & V\_max = 192$\surd$3 cm$^3$ $\approx$ 332,6 cm$^3$ & V\_max = 192$\surd$3 cm$^3$ ($\approx$ 332,55 cm$^3$) & igual & ambos \\
b & altura 4$\surd$3 cm; base quadrada de lado 4$\surd$3 cm (cubo de aresta 4$\surd$3 cm) & cubo: altura e lados da base iguais a 4$\surd$3 cm & igual & ambos \\
\bottomrule\end{longtable}%
\begin{itemize}%
\item%
\textbf{a)} Os dois valores coincidem. Refiz a conta: com d = distância do centro ao plano da face, r$^2$ = 36 - d$^2$, S $\leq$ 2r$^2$ e h $\leq$ 2d. Então V $\leq$ 4(36d - d$^3$), cujo máximo fica em d = 2$\surd$3, o que dá 4$\cdot$2$\surd$3$\cdot$24 = 192$\surd$3 $\approx$ 332,55 cm$^3$. Uma busca numérica aleatória chegou ao mesmo valor. O gabarito usa MA $\geq$ MG sobre a$^2$, b$^2$ e h$^2$ e obtém 48\^{}\{3/2\} = 192$\surd$3.%
\item%
\textbf{b)} Com d = 2$\surd$3 temos r$^2$ = 24, h = 2d = 4$\surd$3 e lado = r$\surd$2 = $\surd$48 = 4$\surd$3. A diagonal do cubo vale 4$\surd$3$\cdot$$\surd$3 = 12 = 2R, então os 8 vértices estão na esfera. Isso é coerente com o gabarito: a$^2$ = b$^2$ = h$^2$ = 48.%
\end{itemize}%
\paragraph{Pontos de dúvida}%
\begin{itemize}%
\item%
\textit{(key)} O gabarito parte de uma base retangular a$\times$b sem justificar, e o enunciado só diz 'quadrangular'. Também supõe implicitamente que os 8 vértices estão na esfera (h = $\surd$(144 - a$^2$ - b$^2$)), embora o enunciado só exija 4. O resultado não muda, porque o máximo geral é atingido pelo cubo. Mesmo assim, como argumento, a resolução do gabarito fica incompleta.%
\item%
\textit{(key)} Erro de digitação na fonte do gabarito: aparece 'a = 4$\surd$(3 cm)' no lugar de 'a = 4$\surd$3 cm'. O sentido não muda.%
\item%
\textit{(solution)} A linha de conferência traz um '$\cdot$\ldots{}' solto em '2$\surd$3/(3$\pi$)$\cdot$\ldots{}$\approx$37\%'. A razão correta é exatamente 2$\surd$3/(3$\pi$) $\approx$ 0,3676. É só um defeito de redação, sem efeito no resultado.%
\end{itemize}%
\end{document}